% =============================================================================
% Chapter: Implementation
% =============================================================================
\chapter{Implementation}
\label{ch:implementation}

This chapter describes the design and integration of Microscaling (MX) format
support into the RedMulE GEMM accelerator. The core challenge is that the
existing systolic array and its surrounding memory subsystem operate exclusively
on FP16 data, while MX-encoded tensors arrive as 8-bit mantissas with shared
exponents packed into 1024-bit TCDM beats. The implementation therefore spans
three areas:

\begin{enumerate}
    \item \textbf{MX datapath architecture}
          (\autoref{sec:mx-codec}--\autoref{sec:w-buffer-scm}): a
          decode-before-compute input pipeline that converts MX blocks to FP16 on
          the fly, and a symmetric encode-after-compute output stage. The codec
          units are designed as format-agnostic modules, parameterized to support
          the full MX element type family (E4M3, E5M2, E3M2, E2M3, E2M1).
    \item \textbf{Tiling extensions}
          (\autoref{sec:tiling}): adaptations to the
          scheduler, memory scheduler, and exponent buffer that enable
          multi-tile execution along all three GEMM dimensions (M, N, K) in MX
          mode.
    \item \textbf{Multi-format support and verification}
          (\autoref{sec:multi-format}--\autoref{sec:verification}): a
          format-selector mechanism for switching between MX element types at
          runtime, a bit-accurate Python golden model, and the debug
          infrastructure used throughout integration.
\end{enumerate}

% TODO: Add top-level block diagram here showing the MX datapath inserted between
% the TCDM/streamer interface and the existing RedMulE engine.


% =============================================================================
% MX Datapath Architecture
% =============================================================================
\section{MX Codec Design}
\label{sec:mx-codec}

% ---- Decoder ----
\subsection{Decoder}
\label{ssec:mx-decoder}

The MX decoder converts a block of MX-encoded FP8 mantissas and their shared
exponent into FP16 values. Each MX block consists of 32 eight-bit elements
sharing a single eight-bit scale exponent. Decoding is a two-step process:
format conversion (FP8$\to$FP16) and exponent scaling.

The decoder is implemented as a two-stage elastic pipeline:

\begin{itemize}
    \item \textbf{Stage~0 (S0):} Latches the input FP8 block and shared
          exponent. In parallel, all 32 elements are combinationally converted
          from FP8 to FP16. For E4M3, this involves rebasing the 4-bit exponent
          from bias~7 to bias~15 and zero-extending the 3-bit mantissa to 10
          bits. E5M2 follows the analogous mapping. Subnormals and special values
          (NaN, Inf) are handled explicitly.
    \item \textbf{Stage~1 (S1):} Applies the MX shared exponent as an additive
          correction $\delta = \text{scale} - 127$ to each element's FP16
          exponent. The result is registered and presented at the output.
          Overflow saturates to FP16 max; underflow flushes to zero.
\end{itemize}

Both stages use elastic handshaking with bubble collapsing: a stage advances
when it is empty or its downstream consumer is ready, yielding a sustained
throughput of one MX block (32 FP16 values) per cycle.

% TODO: Insert pipeline timing diagram (S0/S1 stages, showing throughput).

The conversion logic is implemented as purely combinational functions
(\texttt{fp8\_e4m3\_to\_fp16}, \texttt{fp8\_e5m2\_to\_fp16}) selected by the
\texttt{mx\_format\_i} signal, which makes adding new MX element types
(E3M2, E2M1) a matter of adding the corresponding conversion function and
bias constants.

% TODO: Note the evolution from scalar -> parallel -> pipelined versions
% (commits e15c672a -> 3607c83e -> adbc577d) if relevant for showing the
% design-space exploration.


% ---- Encoder ----
\subsection{Encoder}
\label{ssec:mx-encoder}

The encoder performs the reverse operation: it receives 32 FP16 values from the
engine output and produces an MX-encoded block (32 FP8 mantissas + one shared
exponent). The encoding proceeds through a three-state FSM:

\begin{enumerate}
    \item \textbf{IDLE:} Waits for valid FP16 input.
    \item \textbf{SCAN:} Iterates over groups of input lanes to determine the
          maximum FP16 exponent across the block. This maximum determines the
          shared exponent.
    \item \textbf{ENCODE:} Iterates over the same groups, down-converting each
          FP16 element to FP8 relative to the shared exponent, using
          round-to-nearest-even (RNE). The packed mantissa block and shared
          exponent are output on separate valid/ready channels.
\end{enumerate}

The encoder operates on \texttt{NUM\_LANES} elements per cycle (parameterized,
currently 8 for the 32-wide engine output), requiring $32/\text{NUM\_LANES}$
cycles for the scan pass and equally many for the encode pass. The resulting
throughput is lower than the decoder, but this is acceptable because the engine
produces output at a much lower rate than it consumes input (one Z row every
$N_\text{unpacked}$ engine cycles).

% TODO: Consider adding a table comparing decoder vs. encoder throughput.


% -----------------------------------------------------------------------------
\section{Input Datapath}
\label{sec:mx-input-datapath}
% -----------------------------------------------------------------------------

The MX input datapath is inserted between the HCI streamer interface and the
existing X/W data FIFOs. It consists of four modules connected in series.

% TODO: Insert detailed datapath block diagram.

\subsection{Slot Buffer}
\label{ssec:slot-buffer}

A single 1024-bit TCDM beat contains $1024/8 = 128$ FP8 elements, but the
decoder processes 32 elements at a time (one MX block). The slot buffer
(\texttt{redmule\_mx\_slot\_buffer}) splits each incoming beat into four
256-bit \emph{slots} and buffers them in independent FIFOs for the X and W
paths.

Mantissa slots and exponent data are queued independently, so the exponent
prefetch can proceed even when the mantissa consumer is stalled. For each
stream, the buffer presents a valid slot and its corresponding exponent; a
per-stream \texttt{consume} signal advances to the next entry.

\subsection{Exponent Buffer}
\label{ssec:exp-buffer}

Shared exponents are streamed from TCDM through dedicated source ports
(X\_EXP, W\_EXP) into exponent prefetch buffers
(\texttt{redmule\_exp\_buffer}). Each buffer is a circular SRAM that unpacks
individual exponents from tightly packed beats and stores them sequentially.

The exponent buffer supports three features critical for multi-tile operation:

\begin{itemize}
    \item \textbf{Total count cap:} When the last beat has fewer valid exponents
          than the beat capacity (e.g., 64 valid in a 128-slot beat), the buffer
          discards junk padding based on \texttt{total\_count\_i}.
    \item \textbf{Segment gating:} Limits the number of exponents consumed
          per segment to \texttt{segment\_size\_i}, preventing the decoder from
          reading ahead into future M-tile exponents.
    \item \textbf{Mark/rewind:} Saves and restores the read pointer to replay
          exponents across K-tile and M-tile iterations. Mark is triggered at
          M-tile boundaries (\texttt{x\_rows\_iter\_en}); rewind at K-tile
          boundaries (\texttt{x\_w\_iters\_en}).
\end{itemize}

\subsection{Decoder Scheduling}
\label{ssec:decoder-scheduling}

Two decoder architectures were explored, representing a throughput--complexity
trade-off:

\paragraph{Shared Arbiter (Single Decoder).}
A round-robin arbiter (\texttt{redmule\_mx\_arbiter}) multiplexes X and W slot
pairs into a single decoder instance. The arbiter gives priority to W when the
W FIFO is empty, preventing engine starvation. An elastic buffer (depth~4)
between the decoder output and the W FIFO absorbs bursts. This design is
simple but creates a bottleneck: when the W elastic buffer or W FIFO is full,
the arbiter blocks the decoder entirely, stalling X decoding even when X FIFO
space is available.

\paragraph{Dual Decoder.}
Two independent decoder instances (\texttt{i\_mx\_decoder\_x},
\texttt{i\_mx\_decoder\_w}) each consume directly from their respective slot
buffer outputs. This eliminates cross-path blocking entirely: X and W decode at
full speed regardless of the other path's FIFO state. The input mux target
gating signals (\texttt{target\_is\_x\_i}, \texttt{target\_is\_w\_i}) are tied
high since each decoder always feeds its own path.

The dual decoder provides higher sustained throughput but removes the natural
stall cycles that the arbiter's blocking behavior created at tile boundaries.
This exposed a latent W buffer issue where stale data from the previous K-tile
remained in the SCM (Standard Cell Memory) during the first cycles of a new
tile, which is discussed in \autoref{sec:w-buffer-scm}. A quantitative
comparison of both architectures is presented in \autoref{ch:results}.

\subsection{Input Mux and Beat Packing}
\label{ssec:input-mux}

The input mux (\texttt{redmule\_mx\_input\_mux}) reassembles decoded FP16
values back into 512-bit beats compatible with the downstream data FIFOs. Since
the decoder outputs 32 FP16 values ($32 \times 16 = 512$ bits) per cycle, and
the data FIFO expects 512-bit beats, the packing ratio is $512/512 = 1$
(no multi-cycle accumulation needed at 512-bit FIFO width).

For the W path, the mux additionally tracks the row-chunk position
(\texttt{w\_dec\_pos\_q}) to correctly handle K-tile boundaries where the number
of chunks per W row may vary.

In FP16 bypass mode (\texttt{mx\_enable\_i = 0}), the mux passes the raw
streamer data through unchanged.


% -----------------------------------------------------------------------------
\section{Output Datapath}
\label{sec:mx-output-datapath}
% -----------------------------------------------------------------------------

The output stage (\texttt{redmule\_mx\_output\_stage}) sits between the engine
Z output and the Z stream sink. In MX mode, the FP16 engine output is buffered
in a small FIFO and then fed to the MX encoder, which produces packed FP8
mantissa beats and a separate exponent stream. The mantissa data is forwarded to
the Z FIFO for TCDM write-back; the exponent is written through a dedicated
exponent stream port.

In FP16 mode, the output stage is bypassed and the engine output passes directly
to the Z FIFO.

% TODO: The serialization factor (32 FP16 -> 32 FP8 + 1 exp takes multiple
% cycles) and how the scheduler Z drain logic accounts for this.


% -----------------------------------------------------------------------------
\section{W Buffer}
\label{sec:w-buffer-scm}
% -----------------------------------------------------------------------------

The W buffer (\texttt{redmule\_w\_buffer}) is a shift-register backed by an SCM
(Standard Cell Memory) array of $\text{ARRAY\_HEIGHT} = 32$ rows. Each
\texttt{LOAD\_W} cycle writes one new W row into the SCM and shifts all rows
down by one position. The engine reads \emph{all 32 rows simultaneously} each
cycle (one element per row feeds one row of the systolic array).

This design creates a fill latency: after a K-tile or M-tile boundary, the
first 32 \texttt{LOAD\_W} cycles produce partially stale output. Rows that have
not yet been overwritten still contain data from the previous tile. Under the
shared arbiter design, the natural stalls from W FIFO backpressure provided
enough fill time to mask this effect. With the dual decoder, the faster W
decoding pipeline removes these stalls, exposing the stale data to the engine.

% TODO: Add timing diagram showing the fill progression and stale window.

This is an area of active investigation. Several mitigation approaches were
explored, including SCM data clearing, per-row valid-bit masking, and W row
counter resynchronization. Each approach faced the fundamental constraint that
the engine reads all 32 SCM rows every cycle while only one row is written per
cycle, and the shift operation must continue during the Z buffer fill phase
(\texttt{z\_avail}). The currently most promising solutions are double-buffering
the SCM or introducing a dedicated preload state at M-tile boundaries.

% TODO: Update this section once the fix is finalized.


% =============================================================================
% Tiling and Verification
% =============================================================================
\section{Tiling Extensions}
\label{sec:tiling}

The baseline RedMulE engine supports tiled GEMM execution for matrices larger
than the hardware tile size. For a computation
$Z[M \times K] = X[M \times N] \times W[N \times K] + Y[M \times K]$, the
tiler decomposes the problem into iterations along three dimensions:

\begin{itemize}
    \item \textbf{M-tiling:} \texttt{x\_rows\_iter} $= M / \text{ARRAY\_WIDTH}$
          iterations over groups of 32 output rows.
    \item \textbf{N-tiling:} \texttt{x\_cols\_iter} $= N / \text{TILE}$
          iterations over the reduction dimension (TILE = 64 FP16 elements per
          beat).
    \item \textbf{K-tiling:} \texttt{w\_cols\_iter} $= K / \text{TILE}$
          iterations over groups of 64 output columns.
\end{itemize}

In MX mode, the tiler adjusts the reduction dimension by the MX pack factor:
$N_\text{unpacked} = N_\text{cfg} \times 2$, since each MX FP8 element occupies
half the bus width of an FP16 element. The following sections describe the bugs
encountered when enabling each tiling dimension for MX and the fixes applied.


% -----------------------------------------------------------------------------
\subsection{M-Tiling}
\label{sec:m-tiling}
% -----------------------------------------------------------------------------

M-tiling iterates over groups of \texttt{ARRAY\_WIDTH}~=~32 output rows. Each
M-tile requires a full pass over all N-tiles and K-tiles before advancing to the
next group of rows. Two interacting bugs prevented correct M-tiling in MX mode.

\subsubsection{Z Buffer Fill/Drain Overlap}
\label{ssec:z-drain-overlap}

Between M-tiles, the Z buffer transitions from the drain phase (writing
accumulated results to TCDM) to the fill phase (loading Y-bias for the next
tile). When the fill phase for the next K-tile began while the Z buffer was
still draining, the \texttt{fill\_shift} counter incremented prematurely. Upon
drain completion, the non-zero \texttt{fill\_shift} permanently disabled the
load enable, deadlocking the pipeline.

The fix introduces a \texttt{z\_drain\_guard} signal: a one-cycle-delayed copy
of the Z buffer's \texttt{z\_valid} (drain-active) flag. During the guard
window, the engine clock is gated, preventing the fill counter from advancing
until the drain has fully completed and its effects have propagated. The Z
buffer's \texttt{fill\_shift} and \texttt{d\_index} counters are explicitly
reset on the \texttt{PUSHED}$\to$\texttt{EMPTY} transition.


\subsubsection{Fill/Engine Clock Gate Desynchronization}
\label{ssec:fill-engine-desync}

The Z buffer fill signal (\texttt{cntrl\_z\_buffer\_o.fill}) is combinational,
responding immediately to changes in \texttt{stall\_engine}. The engine clock
gate (\texttt{row\_clk\_en\_q}), however, is registered with a one-cycle delay
through a flip-flop and \texttt{tc\_clk\_gating} latch. When
\texttt{stall\_engine} deasserts during the fill phase:

\begin{enumerate}
    \item The fill counter resumes \emph{immediately} (combinational path).
    \item The engine clock resumes \emph{next cycle} (registered
          \texttt{row\_clk\_en\_q}).
    \item For one cycle, the fill counter captures a value that the engine has
          not yet produced---a duplicate of the previous cycle's output.
\end{enumerate}

This manifested as exactly one duplicated value at fill position 57 in the first
M-tile of a 96$\times$64$\times$64 MX computation (19 errors, all $\pm$1--2 ULP).

The fix aligns fill advancement with the engine clock by gating the fill signal
and \texttt{z\_avail\_counter} with \texttt{row\_clk\_en\_q[0]} (the registered
engine clock enable state) rather than the combinational stall conditions.

% TODO: Consider a timing waveform figure showing the 1-cycle misalignment.


% -----------------------------------------------------------------------------
\subsection{K-Tiling}
\label{sec:k-tiling}
% -----------------------------------------------------------------------------

K-tiling iterates over groups of $\text{TILE} = 64$ output columns. Each K-tile
pass reuses the same X data but loads a new horizontal strip of W. Three issues
arose in the exponent buffer management for K-tiling.

\subsubsection{Exponent Beat Padding}

Each 1024-bit TCDM beat carries up to 128 X exponents (8 bits each). For
non-power-of-two configurations, the last beat may contain fewer than 128 valid
exponents. Without a cap, the buffer stored junk padding values as real
exponents, corrupting the scale factors for subsequent blocks.

The fix adds a \texttt{total\_count\_i} parameter to the exponent buffer,
computed by the memory scheduler from the actual number of exponents for the
current dimension. The buffer discards entries beyond this count when unpacking
the final beat.

\subsubsection{Mark/Rewind Timing}

The exponent buffer's mark and rewind operations were originally triggered by
engine-side signals (\texttt{start\_computation}, \texttt{w\_mat\_iters\_en}).
However, the decoder runs ahead of the engine by several cycles (due to FIFO
buffering), so engine-side triggers caused the buffer to mark or rewind at
positions that were already stale from the decoder's perspective.

The fix retargets both operations to X-side streamer signals:
\begin{itemize}
    \item \textbf{Rewind:} Triggered by \texttt{x\_w\_iters\_en} (K-tile
          boundary on the X memory read side), which aligns with when the decoder
          actually needs to replay exponents.
    \item \textbf{Mark:} Triggered by \texttt{x\_rows\_iter\_en} (M-tile
          boundary on the X memory read side), saving the pointer at the start of
          each M-tile's exponent range.
\end{itemize}

\subsubsection{Segment Gating}

Without segment gating, the decoder consumed exponents as fast as slots were
available, potentially reading into exponents belonging to the next M-tile before
the current one completed. The \texttt{segment\_size\_i} parameter limits the
exponent buffer to emit at most $N_\text{unpacked}$ exponents per segment
(between consecutive mark/rewind events), enforcing M-tile isolation.


% -----------------------------------------------------------------------------
\subsection{N-Tiling}
\label{sec:n-tiling}
% -----------------------------------------------------------------------------

N-tiling iterates over the reduction dimension in groups of
$\text{TILE} = 64$ elements. This was the final tiling dimension to be enabled
and required fixes in both the software test infrastructure and the RTL.

\subsubsection{X Memory Layout Reordering}

The original test vector generator produced X data in N-tile-major order: all M
rows for N-tile~0, then all M rows for N-tile~1, etc. The memory scheduler,
however, reads linearly per M-tile. For multi-M-tile configurations, M-tile~0's
reader would consume M-tile~1's data from N-tile~0 instead of its own N-tile~1
data.

The fix reorders the X memory layout to M-tile-major order: within each M-tile,
all N-tiles are stored contiguously. This was implemented in the Python test
vector generator (\texttt{gen\_mx\_test\_vectors.py}) with a
\texttt{reorder\_mn\_tile\_major()} function parameterized by
\texttt{--m-tile-rows} (= ARRAY\_WIDTH = 32).

\subsubsection{Address Offset Advancement}

The \texttt{x\_rows\_offs\_q} register in the memory scheduler tracks the TCDM
address offset for each M-tile's X data. The update condition compared the
column iteration counter (\texttt{x\_cols\_iters\_q}) against a tiler register
(\texttt{X\_ITERS[15:0]}), which holds the original column iteration count. In
MX mode, however, the effective column iteration limit
(\texttt{msched\_x\_cols\_limit}) differs from the stored value. For $N = 64$
both values coincidentally equal~1, so the bug was invisible. For $N > 64$ (e.g.,
$N = 128 \Rightarrow \texttt{x\_cols\_iter} = 2$), the comparison was never
true, and the address offset never advanced---all M-tiles read from the same
base address.

The fix changes the comparison target to \texttt{msched\_x\_cols\_limit - 1},
consistent with how the row iteration update already operated.


% -----------------------------------------------------------------------------
\subsection{Multi-Format Support}
\label{sec:multi-format}
% -----------------------------------------------------------------------------

The MX specification defines multiple element types with different
exponent/mantissa splits. The current implementation supports E4M3 and E5M2;
E3M2 (FP6) and E2M1 (FP4) are planned extensions.

The format selector is encoded as a 3-bit enum (\texttt{mx\_format\_e}) carried
through the control path from the software configuration register to the decoder
and encoder. The decoder selects the appropriate conversion function (bias
constant, mantissa width, special-value handling) based on this signal. The
encoder symmetrically adjusts its down-conversion logic.

Adding E3M2 and E2M1 requires:
\begin{itemize}
    \item New conversion functions in the decoder and encoder.
    \item Adjusted pack factors: FP6 elements are 6 bits wide
          ($\lfloor 8/6 \rfloor$ packing requires padding handling), and FP4
          elements are 4 bits wide ($256/4 = 64$ elements per 256-bit slot).
    \item Updated slot buffer unpacking logic to handle non-byte-aligned element
          widths (FP6) or increased element counts (FP4).
    \item Golden model extensions for bit-accurate reference computation.
\end{itemize}

% TODO: Expand this section after FP4/FP6 implementation is complete.


% -----------------------------------------------------------------------------
\subsection{Verification}
\label{sec:verification}
% -----------------------------------------------------------------------------

\subsubsection{Golden Model}
\label{ssec:golden-model}

A bit-accurate Python golden model (\texttt{mx\_fp\_golden.py}) serves as the
reference for simulation comparison. The model performs the full MX GEMM
pipeline: encodes FP32 input matrices to MX format (shared exponent selection,
FP8 quantization with RNE), decodes back to FP16, computes the GEMM in FP16
arithmetic, and encodes the result. This mirrors the hardware pipeline exactly,
ensuring that any numerical difference is a true hardware bug rather than a
precision artifact.

The test vector generator (\texttt{gen\_mx\_test\_vectors.py}) produces
TCDM-layout binary stimuli with the correct M-tile-major ordering, exponent
packing, and byte alignment. A companion script
(\texttt{gen\_mx\_golden.py}) generates the expected output in the same binary
format for comparison.

\subsubsection{Debug Methodology}
\label{ssec:debug-methodology}

Several techniques were developed during integration to diagnose the tiling bugs
described in \autoref{sec:m-tiling}--\autoref{sec:n-tiling}:

\begin{itemize}
    \item \textbf{W increment override (\texttt{+W\_INCR\_OVERRIDE}):} A
          simulation-time flag that replaces the W buffer output with
          monotonically incrementing FP16 values. This isolates W-data-dependent
          errors from engine pipeline issues by making the W contribution
          deterministic and position-traceable.
    \item \textbf{Per-stage display dumps:} \texttt{\$display}-based probes at
          the slot buffer output, decoder output, input mux, and engine input
          that log cycle-by-cycle data flow, enabling visual diff against the
          golden model.
    \item \textbf{Engine input probes:} Direct observation of the values
          presented to the systolic array on each engine cycle, with M-tile,
          K-tile, and W-row annotations. This was instrumental in identifying
          the W buffer SCM stale data issue.
\end{itemize}

The full test matrix, cycle counts, and correctness results are presented in
\autoref{ch:results}.
